\section{Práctica de laboratorio}

La práctica de laboratorio se dividió en dos etapas: ensayos de seguridad y ensayos de eficacia de equipos médicos. En ambas etapas se utilizaron analizadores para realizar las mediciones correspondientes y comparar los resultados obtenidos con los valores establecidos por la normativa y/o los manuales de usuario de los equipos bajo prueba.

\subsection{Ensayos de seguridad}

El objetivo de esta etapa fue evaluar diferentes parámetros relacionados con la seguridad eléctrica de los equipos médicos. Las mediciones realizadas fueron la resistencia a tierra, el consumo de corriente, las corrientes de fuga y la resistencia de aislamiento.

\subsubsection{Resistencia a tierra}

La prueba de resistencia de la línea de tierra se realizó con el equipo médico apagado. Esta medición permite determinar la impedancia entre el terminal de tierra de protección (\textit{PE}) del analizador y las partes conductoras expuestas del equipo que se encuentran conectadas a la tierra de protección.

Antes de realizar la medición se efectuó la calibración del analizador de forma que el punto de referencia sea cero, eliminando la medición de la resistencia asociada al cable utilizado.

Una vez finalizada la medición de resistencia a tierra, las pruebas restantes se realizaron con el equipo médico encendido. 

\subsubsection{Consumo de corriente}

La prueba de consumo de corriente permitió determinar la corriente consumida por el dispositivo bajo prueba a través de las conexiones de red del receptáculo de prueba del analizador.

\subsubsection{Corrientes de fuga}

Se realizaron mediciones de corrientes de fuga para diferentes configuraciones del equipo bajo prueba. Entre estas se encuentran la corriente de fuga a tierra, la de la carcasa o por contacto, la de las partes aplicadas y la asociada a fase y tierra de protección.

\subsubsection{Resistencia de aislamiento}

Finalmente, se realizaron pruebas de resistencia de aislamiento entre la red y la tierra de protección, las partes aplicadas y la tierra de protección, y la red y las partes aplicadas. Los ensayos se realizaron utilizando una tensión continua de 500 V.

\subsection{Ensayos de eficacia}

En la segunda etapa se evaluó la eficacia de los equipos médicos mediante la comparación de sus mediciones con los valores generados por los analizadores correspondientes.

Para esta etapa se utilizaron el FLUKE ProSim SPOT Light y el FLUKE ProSim4. Se ensayaron un monitor multiparamétrico y un oxímetro de pulso.

Los valores obtenidos durante las pruebas fueron registrados para posteriormente compararlos con los valores de referencia proporcionados por los analizadores y evaluar el funcionamiento de los equipos médicos.